% Generated draft: requires mathematical and source review.
\documentclass[11pt]{article}
\usepackage[margin=25mm]{geometry}
\usepackage{fontspec}
\setmainfont{Latin Modern Roman}
\usepackage{amsmath,amssymb,mathtools,unicode-math}
\setmathfont{Latin Modern Math}
\usepackage{enumitem,needspace,xurl,hyperref}
\setlength{\parindent}{0pt}
\setlength{\parskip}{5pt}
\setlength{\emergencystretch}{3em}
\newcommand{\N}{\mathbb{N}}
\newcommand{\Z}{\mathbb{Z}}
\newcommand{\Q}{\mathbb{Q}}
\newcommand{\R}{\mathbb{R}}
\newcommand{\C}{\mathbb{C}}
\newcommand{\F}{\mathbb{F}}
\newcommand{\A}{\mathbb{A}}
\newcommand{\PP}{\mathbb{P}}
\newcommand{\E}{\mathbb{E}}
\newcommand{\eps}{\varepsilon}
\newcommand{\dd}{\,\mathrm{d}}
\begin{document}
\section*{UnsolvedMath ID: 20002372}
\textbf{Source problem code:} catalogue entry
\subsection{Definitions and mathematical statement}
Let $m,n\geq1$, and let $f_1,\ldots,f_n\in\C[x_1,\ldots,x_m]$ be homogeneous of the same positive degree $d$, with
\[
V(f_1,\ldots,f_n)=\{(0,\ldots,0)\}.
\]
Question 36 (Jarden) asks to prove that for every $\bar b=(b_1,\ldots,b_n)\in\C^n$ the affine fiber
\[
\{\bar x\in\C^m:f_i(\bar x)=b_i\text{ for }1\leq i\leq n\}
\]
is finite, including the possibility of an empty fiber. Positive degree and positive numbers of variables and forms are the conventional nontrivial setting; the source does not explicitly list these integer bounds.

The source already includes its qualitative solution: homogenize to $f_i(x_1,\ldots,x_m)-b_i x_0^d=0$ in $\PP^m$. An infinite affine fiber would yield a positive-dimensional projective closure meeting $x_0=0$, giving a common nonzero projective zero of the homogeneous forms, contrary to the hypothesis. This supplied solution is part of the record, not a newly proposed solution.

The aid states a stronger characteristic-free refinement. For any algebraically closed field $k$ and forms satisfying the analogous hypotheses, put $R=k[x_1,\ldots,x_m]$, $I=(f_1,\ldots,f_n)$, and
\[
F:\A_k^m\longrightarrow\A_k^n,\qquad\bar x\longmapsto(f_1(\bar x),\ldots,f_n(\bar x)).
\]
Then $F$ is a finite morphism, and, for every $\bar b\in k^n$,
\[
\operatorname{length}_k F^{-1}(\bar b)
=\dim_k R/(f_1-b_1,\ldots,f_n-b_n)
\leq\dim_k R/I\leq d^m.
\]
Thus the number of distinct points in every fiber is at most $d^m$, while the length counts multiplicity. Finiteness of the morphism does not assert surjectivity onto $\A_k^n$. The aid says the scheme-length bound is sharp in every characteristic, as illustrated by $n=m$, $f_i=x_i^d$ and the fiber over zero; it says the distinct-point bound is sharp when $\operatorname{char}(k)\nmid d$, using a fiber with all coordinates nonzero in the same example. No claim of reduced fibers in other characteristics is made.
\subsection{Short English statement}
Homogeneous forms with only the origin as a common zero have finite fibers over every target value. The source includes a projective proof. The supplied refinement makes the polynomial map finite and bounds every fiber's scheme length, hence its number of points, by $d^m$, with the stated sharpness qualifications.
\subsection{Sources}
\begin{itemize}
\item Local UnsolvedMath registry, permanent ID $20002372$, \texttt{catalogue entry}; \texttt{aim-logic-notes.json}, \texttt{notes[147]}.
\item American Institute of Mathematics, \emph{Extensions of Hilbert's tenth problem}, Question 36 (Jarden), including the embedded projective-homogenization solution: \url{https://aimath.org/WWN/hilberts10th/hilberts10th.pdf}. Supplied verification reports correction of subscripts, vector notation, and the run-together ``C.Solution'' extraction.
\item Research attempt 1: full solution, problem solved at run, rigor label ``proved,'' novelty confidence low. Registry classification solved in literature, status solved, statement status corrected/verified. No separate authoritative reference is supplied for the quantitative scheme-length refinement.
\end{itemize}
\end{document}
